\begin{center}\Large\bfseries MỤC LỤC\end{center}
\vspace{0.3cm}
\begin{spacing}{1.08}
\begin{longtable}{@{}p{0.86\textwidth}r@{}}
\textbf{LỜI CẢM ƠN} & \textbf{i}\\
\textbf{LỜI CAM ĐOAN} & \textbf{ii}\\
\textbf{DANH MỤC TỪ VIẾT TẮT} & \textbf{vii}\\
\textbf{MỞ ĐẦU} & \textbf{1}\\
\textbf{CHƯƠNG 1: CƠ SỞ LÝ THUYẾT} & \textbf{4}\\
\hspace*{0.35cm} 1.1. Chatbot trong hỗ trợ tra cứu thủ tục hành chính & 4\\
\hspace*{0.35cm} 1.2. Xử lý ngôn ngữ tự nhiên và mô hình ngôn ngữ lớn & 4\\
\hspace*{0.35cm} 1.3. Embedding, tìm kiếm ngữ nghĩa và reranking & 5\\
\hspace*{0.35cm} 1.4. Kỹ thuật Retrieval-Augmented Generation (RAG) & 6\\
\hspace*{0.35cm} 1.5. Vector database và ChromaDB & 7\\
\hspace*{0.35cm} 1.6. Đánh giá chất lượng hệ thống RAG & 8\\
\hspace*{0.35cm} 1.7. Công nghệ web và lưu trữ được áp dụng với hệ thống & 9\\
\textbf{CHƯƠNG 2: KHẢO SÁT VÀ PHÂN TÍCH HỆ THỐNG} & \textbf{10}\\
\hspace*{0.35cm} 2.1. Khảo sát thực trạng tra cứu thủ tục hành chính & 10\\
\hspace*{0.35cm} 2.2. Phân tích dữ liệu đầu vào & 12\\
\hspace*{0.35cm} 2.3. Phân tích bài toán và phạm vi hệ thống & 14\\
\hspace*{0.35cm} 2.4. Phân tích tác nhân và nhu cầu sử dụng & 15\\
\hspace*{0.35cm} 2.5. Phân tích quy trình nghiệp vụ & 16\\
\hspace*{0.75cm} 2.5.1. Quy trình người dân tra cứu thủ tục hành chính & 16\\
\hspace*{0.75cm} 2.5.2. Quy trình hỏi đáp thủ tục hành chính bằng chatbot & 16\\
\hspace*{0.75cm} 2.5.3. Quy trình quản trị dữ liệu thủ tục hành chính & 17\\
\hspace*{0.35cm} 2.6. Phân tích yêu cầu hệ thống & 18\\
\hspace*{0.75cm} 2.6.1. Yêu cầu chức năng & 18\\
\hspace*{0.75cm} 2.6.2. Yêu cầu phi chức năng và tiêu chí chấp nhận & 19\\
\hspace*{0.35cm} 2.7. Sơ đồ Use Case và đặc tả Use Case chính & 20\\
\hspace*{0.75cm} 2.7.1. Sơ đồ Use Case tổng quát & 21\\
\hspace*{0.75cm} 2.7.2. Sơ đồ Use Case nhóm người dùng & 22\\
\hspace*{0.75cm} 2.7.3. Sơ đồ Use Case nhóm quản trị viên & 22\\
\hspace*{0.75cm} 2.7.4. Đặc tả các Use Case chính & 23\\
\hspace*{0.35cm} 2.8. Tổng hợp phân tích và lựa chọn giải pháp RAG & 28\\
\textbf{CHƯƠNG 3: THIẾT KẾ HỆ THỐNG} & \textbf{30}\\
\hspace*{0.35cm} 3.1. Thiết kế kiến trúc tổng thể & 30\\
\hspace*{0.35cm} 3.2. Thiết kế mô hình lớp thực thể và quan hệ & 31\\
\hspace*{0.35cm} 3.3. Thiết kế cơ sở dữ liệu & 34\\
\hspace*{0.35cm} 3.4. Thiết kế pipeline RAG và trích dẫn nguồn & 36\\
\hspace*{0.35cm} 3.5. Thiết kế API và request/response mẫu & 40\\
\hspace*{0.35cm} 3.6. Thiết kế giao diện người dùng và quản trị & 46\\
\hspace*{0.35cm} 3.7. Thiết kế triển khai và cấu trúc mã nguồn & 48\\
\textbf{CHƯƠNG 4: TRIỂN KHAI} & \textbf{54}\\
\hspace*{0.35cm} 4.1. Cài đặt backend và pipeline RAG & 54\\
\hspace*{0.35cm} 4.2. Cài đặt frontend và kết nối API & 54\\
\hspace*{0.35cm} 4.3. Bảo vệ hệ thống và dữ liệu & 55\\
\hspace*{0.35cm} 4.4. Hoạt động thử nghiệm, kiểm thử và đánh giá định lượng & 57\\
\hspace*{0.35cm} 4.5. Đánh giá kết quả triển khai và hướng dẫn sử dụng & 66\\
\textbf{KẾT LUẬN VÀ HƯỚNG PHÁT TRIỂN} & \textbf{69}\\
\textbf{TÀI LIỆU THAM KHẢO} & \textbf{70}\\
\textbf{PHỤ LỤC} & \textbf{74}\\
\textbf{NHẬN XÉT CỦA GIÁO VIÊN HƯỚNG DẪN} & \textbf{84}\\
\end{longtable}
\end{spacing}